\documentclass[11pt]{amsart}

\usepackage[T1]{fontenc}
\usepackage{lmodern}
\usepackage{microtype}
\usepackage{mathtools,amssymb,amsthm}
\usepackage[hidelinks]{hyperref}

\hypersetup{
  pdftitle={The face-angle map of a tetrahedron is not injective on the danger cylinder}
}

\newtheorem{theorem}{Theorem}
\newtheorem{lemma}[theorem]{Lemma}
\theoremstyle{remark}
\newtheorem*{problem}{Problem 3 \textup{(Nikitenko--Nikonorov)}}

\DeclareMathOperator{\Cyl}{Cyl}

\title[The face-angle map is not injective on the danger cylinder]
{The face-angle map of a tetrahedron is not injective on the danger cylinder}

\date{}

\subjclass[2020]{51M04, 51M20, 68T45}
\keywords{face angle, tetrahedron with a given base, danger cylinder, P3P problem,
self-intersection, counterexample}

\begin{document}

\begin{abstract}
Fix a nondegenerate triangle $ABC$ in the plane $z=0$ with circumcircle $S$, let
$\Cyl=S\times\mathbb R$, let $\Pi^{+}=\{(p,q,r):r>0\}$, and let
$F(D)=(\cos\angle BDC,\cos\angle CDA,\cos\angle ADB)$ be the triple of face angles at the
apex $D$ of the tetrahedron $ABCD$. Problem~3 of Nikitenko and Nikonorov asks whether $F$
is injective on $\Cyl\cap\Pi^{+}$, stated there in the equivalent form that the surface
$F(\Cyl\cap\Pi^{+})$ has no self-intersection. The answer is no. For the base
$A=(1,0,0)$, $B=(\tfrac{15}{17},\tfrac{8}{17},0)$, $C=(0,-1,0)$ on the unit circle and the
base point $P=(\tfrac45,\tfrac35,0)$, the two apexes $D_j=(\tfrac45,\tfrac35,r_j)$ whose
heights satisfy $r_1^2,r_2^2=(80\mp12\sqrt{35})/425$ are distinct points of
$\Cyl\cap\Pi^{+}$ with
\[
 F(D_1)=F(D_2)=\Bigl(\tfrac{2}{51}\sqrt{119},\ \tfrac{2}{13}\sqrt{21},\
 \tfrac{2}{17}\sqrt{51}\Bigr).
\]
All the coordinates defining the base and the base point are rational, and the
verification reduces to two rational identities per coordinate.
\end{abstract}

\maketitle

\section{The problem}

Let $ABC$ be a nondegenerate triangle in the plane $z=0$, with circumcircle $S$ of centre
$O$ and radius $R$. For a point $D=(p,q,r)$ with $r\neq0$ the four points $A,B,C,D$ span a
tetrahedron, and Nikitenko and Nikonorov \cite{NN} study the map that records its three
face angles at the apex,
\begin{equation}\label{eq:F}
 F(D)=\bigl(\cos\angle BDC,\ \cos\angle CDA,\ \cos\angle ADB\bigr),
\end{equation}
on the half-space $\Pi^{+}=\{(p,q,r):r>0\}$. Writing $\Cyl=S\times(-\infty,\infty)$ for the
right cylinder over the circumcircle and $FC=F(\Cyl\cap\Pi^{+})$ for its image, they ask
\cite[Problem~3]{NN}:

\begin{problem}
\sloppy
Is it true that the surface $FC=F(\Cyl\cap\Pi^{+})$ has no self-intersection, i.e.\ the map
$F$ is injective on the set $\Cyl\cap\Pi^{+}$?
\end{problem}

\noindent The quoted statement itself identifies the absence of a self-intersection of $FC$
with injectivity of $F$ on $\Cyl\cap\Pi^{+}$, and it is that injectivity which we refute.
Its negation is existential, so a single explicit configuration decides it; we exhibit one
whose base and base point have rational coordinates.

Throughout we place $S$ as the unit circle in the plane $z=0$ centred at the origin, so
$R=1$; this is the normalisation of \cite{NN}. The cylinder $\Cyl$ over the circumcircle is
known in computer vision as the \emph{danger cylinder} of the perspective-three-point
problem.

\section{The counterexample}

\begin{theorem}\label{thm:main}
In the plane $z=0$ let
\[
 A=(1,0,0),\qquad B=\Bigl(\tfrac{15}{17},\tfrac{8}{17},0\Bigr),\qquad C=(0,-1,0),
\]
so that $A,B,C$ are rational points of the unit circle $S$ and $ABC$ is nondegenerate, and
let
\[
 P=\Bigl(\tfrac45,\tfrac35,0\Bigr)\in S .
\]
Let $s_1<s_2$ be the two roots of
\begin{equation}\label{eq:quad}
 s^{2}-\tfrac{32}{85}\,s+\tfrac{16}{2125}=0,
 \qquad\text{that is}\qquad
 s_{1,2}=\frac{80\mp12\sqrt{35}}{425},
\end{equation}
and put $r_j=\sqrt{s_j}$ and $D_j=\bigl(\tfrac45,\tfrac35,r_j\bigr)$ for $j=1,2$. Then
$D_1$ and $D_2$ are two distinct points of $\Cyl\cap\Pi^{+}$ and
\begin{equation}\label{eq:image}
 F(D_1)=F(D_2)
 =\Bigl(\tfrac{2}{51}\sqrt{119},\ \tfrac{2}{13}\sqrt{21},\ \tfrac{2}{17}\sqrt{51}\Bigr).
\end{equation}
Consequently $F$ is not injective on $\Cyl\cap\Pi^{+}$, and Problem~3 is answered in the
negative.
\end{theorem}

The two tetrahedra $ABCD_1$ and $ABCD_2$ share the base $ABC$ and all three face angles at
the apex, and their apexes differ: $|D_1A|^{2}=\tfrac25+s_1\neq\tfrac25+s_2=|D_2A|^{2}$.

\section{Proof}

\subsection*{The reduction along one generator}
Fix $P\in S$ and let $D=(P,r)$ run up the vertical generator of $\Cyl$ through $P$. Put
\[
 s=r^{2},\qquad u=|PB|^{2},\quad v=|PC|^{2},\quad w=|PA|^{2},
\]
\[
 a=|BC|,\qquad b=|CA|,\qquad c=|AB|,
\]
and
\[
 K_1=\tfrac12(u+v-a^{2}),\qquad K_2=\tfrac12(v+w-b^{2}),\qquad K_3=\tfrac12(w+u-c^{2}).
\]
Because $D$ lies vertically above $P$ we have $|DB|^{2}=u+s$, $|DC|^{2}=v+s$ and
$|DA|^{2}=w+s$, so the law of cosines in the triangle $BDC$ gives
\begin{equation}\label{eq:g}
 \cos\angle BDC=\frac{K_1+s}{\sqrt{(u+s)(v+s)}},
\end{equation}
and cyclically for the other two coordinates: each coordinate of $F$ along the generator is
a function of the single variable $s$ of the shape
$g(s)=(K+s)\bigl/\sqrt{(p+s)(q+s)}$.

\subsection*{Equal values at two heights}
Suppose $s_1\neq s_2$ are positive and put $\sigma=s_1+s_2$, $\pi=s_1s_2$, so that $s_1,s_2$
are the roots of $s^{2}-\sigma s+\pi$. Reducing modulo that quadratic replaces $s^{2}$ by
$\sigma s-\pi$, and therefore, for a rational constant $\gamma$,
\[
 (K+s)^{2}-\gamma\,(p+s)(q+s)
 \equiv
 \bigl[(K^{2}-\pi)-\gamma\,(pq-\pi)\bigr]
 +\bigl[(2K+\sigma)-\gamma\,(p+q+\sigma)\bigr]s .
\]
The right-hand side is linear in $s$, so it vanishes at both roots if and only if both
brackets vanish. Hence:

\begin{lemma}\label{lem:two}
Let $s_1\neq s_2$ be the roots of $s^{2}-\sigma s+\pi$ and let $p,q,K$ be given. Then
$g(s_1)^{2}=g(s_2)^{2}=\gamma$ holds with
\begin{equation}\label{eq:gamma}
 \gamma=\frac{2K+\sigma}{p+q+\sigma}
 \qquad\text{if and only if}\qquad
 K^{2}-\pi=\gamma\,(pq-\pi),
\end{equation}
provided $p+q+\sigma\neq0$ and $(p+s_j)(q+s_j)\neq0$. If moreover $K+s_1$ and $K+s_2$ have
the same sign, then $g(s_1)=g(s_2)$.
\end{lemma}

\begin{proof}
Immediate from the displayed reduction: the coefficient of $s$ vanishes exactly when
$\gamma$ has the stated value, and then the constant term vanishes exactly when the stated
identity holds. For the last sentence, $g(s_1)^{2}=g(s_2)^{2}$ gives
$g(s_1)=\pm g(s_2)$; the denominators in \eqref{eq:g} are positive, so the sign of $g(s_j)$
is the sign of $K+s_j$, and equal signs force the $+$ case.
\end{proof}

\subsection*{Proof of Theorem \ref{thm:main}}
Every quantity below is a rational number computed from the rational coordinates of
$A,B,C,P$ in Theorem~\ref{thm:main}; a reader can confirm each in a line.

The four points lie on the unit circle, since $15^{2}+8^{2}=17^{2}$ and
$4^{2}+3^{2}=5^{2}$; the base is nondegenerate, three distinct points of a circle never
being collinear; and $P\notin\{A,B,C\}$. The six squared distances are
\begin{equation}\label{eq:lengths}
 \begin{gathered}
 u=|PB|^{2}=\tfrac{2}{85},\qquad v=|PC|^{2}=\tfrac{16}{5},\qquad w=|PA|^{2}=\tfrac25,\\
 a^{2}=\tfrac{50}{17},\qquad b^{2}=2,\qquad c^{2}=\tfrac{4}{17},
 \end{gathered}
\end{equation}
whence
\begin{equation}\label{eq:sigmapi}
 \sigma:=4-(u+v+w)=4-\tfrac{308}{85}=\tfrac{32}{85},
 \qquad
 \pi:=\tfrac14\,uvw=\tfrac{16}{2125},
\end{equation}
and
\[
 K_1=\tfrac{12}{85},\qquad K_2=\tfrac45,\qquad K_3=\tfrac{8}{85}.
\]
The quadratic \eqref{eq:quad} is exactly $s^{2}-\sigma s+\pi$, its discriminant is
\begin{equation}\label{eq:disc}
 \sigma^{2}-4\pi=\Bigl(\tfrac{32}{85}\Bigr)^{2}-\tfrac{64}{2125}
 =\tfrac{4032}{36125}>0,
\end{equation}
and $\sigma>0$, $\pi>0$, so its roots $s_1,s_2$ are two \emph{distinct positive} reals. Hence
$r_1=\sqrt{s_1}$ and $r_2=\sqrt{s_2}$ are distinct positive heights, $D_1\neq D_2$, both
$D_j$ lie in $\Pi^{+}$, and both have horizontal part $P\in S$, so both lie on $\Cyl$: they
are two distinct points of $\Cyl\cap\Pi^{+}$.

Now apply Lemma~\ref{lem:two} once per coordinate, taking $(K,p,q)$ to be $(K_1,u,v)$, then
$(K_2,v,w)$, then $(K_3,w,u)$; all denominators below are nonzero. The values
\eqref{eq:gamma} are
\[
 \gamma_1=\frac{\tfrac{24}{85}+\tfrac{32}{85}}{\tfrac{2}{85}+\tfrac{16}{5}+\tfrac{32}{85}}
 =\frac{56/85}{306/85}=\frac{28}{153},
\]
\[
 \gamma_2=\frac{\tfrac85+\tfrac{32}{85}}{\tfrac{16}{5}+\tfrac25+\tfrac{32}{85}}
 =\frac{168/85}{338/85}=\frac{84}{169},
 \qquad
 \gamma_3=\frac{\tfrac{16}{85}+\tfrac{32}{85}}{\tfrac25+\tfrac{2}{85}+\tfrac{32}{85}}
 =\frac{48/85}{68/85}=\frac{12}{17},
\]
and the three constant-term identities hold:
\begin{align}
 K_1^{2}-\pi&=\tfrac{144}{7225}-\tfrac{16}{2125}=\tfrac{448}{36125}
   =\tfrac{28}{153}\cdot\tfrac{144}{2125}=\gamma_1(uv-\pi),\notag\\
 K_2^{2}-\pi&=\tfrac{16}{25}-\tfrac{16}{2125}=\tfrac{1344}{2125}
   =\tfrac{84}{169}\cdot\tfrac{2704}{2125}=\gamma_2(vw-\pi),\label{eq:three}\\
 K_3^{2}-\pi&=\tfrac{64}{7225}-\tfrac{16}{2125}=\tfrac{48}{36125}
   =\tfrac{12}{17}\cdot\tfrac{4}{2125}=\gamma_3(wu-\pi).\notag
\end{align}
Finally $K_1,K_2,K_3>0$ and $s_1,s_2>0$, so $K_i+s_1$ and $K_i+s_2$ are positive for every
$i$, and Lemma~\ref{lem:two} upgrades the equality of squares to equality of the cosines
themselves. Therefore
\[
 F(D_1)=F(D_2)=\bigl(\sqrt{\gamma_1},\sqrt{\gamma_2},\sqrt{\gamma_3}\bigr)
 =\Bigl(\tfrac{2}{51}\sqrt{119},\ \tfrac{2}{13}\sqrt{21},\ \tfrac{2}{17}\sqrt{51}\Bigr),
\]
which is \eqref{eq:image}, using $\tfrac{28}{153}=\tfrac{4\cdot119}{51^{2}}$,
$\tfrac{84}{169}=\tfrac{4\cdot21}{13^{2}}$ and
$\tfrac{12}{17}=\tfrac{4\cdot51}{17^{2}}$. \qed

\section{Scope}\label{sec:scope}

Theorem~\ref{thm:main} refutes Problem~3 as printed: $F$ is not injective on
$\Cyl\cap\Pi^{+}$. Nothing more is claimed; in particular the non-injectivity locus of $F$
on $\Cyl\cap\Pi^{+}$ is not described here.

\section{Verification}

Theorem~\ref{thm:main} needs no computer: the object is four rational points of a circle, and its proof
is the three rational identities \eqref{eq:three} together with three sign checks. A
program \texttt{verify.py} and a recorded run \texttt{verify.output.txt} accompany this note
nonetheless. The program reads the four points $A,B,C,P$ printed in Theorem~\ref{thm:main}
and nothing else, recomputes \eqref{eq:lengths}--\eqref{eq:disc}, the three $\gamma_i$ and
the three identities \eqref{eq:three}, and then decides $F(D_1)=F(D_2)$ in exact arithmetic
in the real quadratic field $\mathbb Q(\sqrt{35})$ -- including an exact sign test, so no
floating-point number enters that decision. It also evaluates the two face-angle triples by
three-dimensional dot products of the edge vectors in high-precision arithmetic; that route
reports only numerical agreement to the precision used, and the equality of the triples is
established by the exact checks alone. Every check in the recorded run passes.
Lemma~\ref{lem:two} is proved above by hand and is not re-derived by the program.

The recorded run is not confined to this note. Parts of it, and of the scope list it prints,
refer to a Proposition~3, a Remark~3 and Sections~4 and~5 that do not occur here, and to a
rational scan, a control computation, decimal expansions and properties of the base triangle
that support no statement made here; those lines should be read as bearing on other
statements, and nothing in Theorem~\ref{thm:main} depends on them. What the run establishes for this note is the
exhibited object: every number and every inequality of Theorem~\ref{thm:main}, re-derived
from the four points $A,B,C,P$.

\begin{thebibliography}{9}

\bibitem{NN}
E.~V. Nikitenko and Yu.~G. Nikonorov,
\emph{On face angles of tetrahedra with a given base},
Results Math. \textbf{81} (2026), no.~3, Paper No.~60.
\href{https://doi.org/10.1007/s00025-026-02610-x}{doi:10.1007/s00025-026-02610-x};
\href{https://arxiv.org/abs/2505.22374v4}{arXiv:2505.22374v4}.
Problem~3 is quoted from the e-print, version~4.

\end{thebibliography}

\end{document}
